\section{AT4: a direct actual-state comparison and its failed precision target}
\label{sec:at4}

\textbf{Contribution \projecttag{HNM-C-AT4}.} The first contract fixed
$\tau=10^{-8}$, $s=1$, $L=10^4$ and absolute point error $10^{-6}$. Its reviewed
verdict is \emph{limited}: an actual-state interval is obtained, but it fails the
specified precision. The following derivation records the sharper forward
certificate; the reverse proof independently supplied a valid broader tail
bound. The effect-bound refinement was separately checked after both reports
froze and is not claimed as a duplicated independent discovery.

\subsection{Local state error}
At the zero selected triple the dimensionless on-site reference is
$h_b=8\sum_{e\text{ owned by }b}C_e$. Its first positive full-space energy is
$8(3/4)=6$. Therefore, with the Haar product projection $P_R$,
\begin{equation}
 h_R\ge6(I-P_R).
 \label{eq:referencegap}
\end{equation}
Reset the finite interacting ground density on $R$ to $P_R$, preserving its
exterior marginal. Each of the seven incident whole stars changes its
expectation by at most $2(7|\tau|)$ in AQ normalized units. The variational
principle and cancellation of unchanged exterior energy give
\begin{equation}
 \tr(\rho_{N,R}h_R)\le98|\tau|.
 \label{eq:reset}
\end{equation}
Positive spectral truncations justify the cancellation for unbounded reference
operators. The reset trial and finite ground have finite form energy because
the finite interaction is bounded. Pass bounded truncations through the actual
AQ local trace-norm limit, then take their monotone supremum. Thus
\begin{equation}
 1-\tr(\rho_{\tau,R}P_R)\le\frac{49|\tau|}{3},\qquad
 \norm{\rho_{\tau,R}-P_R}_1\le D:=2\sqrt{\frac{49|\tau|}{3}}.
 \label{eq:D}
\end{equation}
The second inequality follows by pure-state trace distance, convexity and
Cauchy--Schwarz. It is uniform for the AQ subsequential states under this
construction; it does not identify them. No convergence-rate assertion is used.

Trace duality gives $|m_\tau|\le D$. A trace-zero density difference of norm at
most $D$ has expectation at most $D/2$ against an effect. Applying this to
$0\le W^2\le I$ gives
\begin{equation}
 \abs{\omega_\tau(W^2)-\tfrac14}\le D/2,
 \qquad v_\tau:=\norm{\chi_\tau}^2\le\tfrac14+D/2.
 \label{eq:effect}
\end{equation}
The complex operator $W\alpha_t^0(W)$ below is not an effect; its error remains
$D$, not $D/2$.

\subsection{Dynamics and centering}
In each finite box split $H_{N,\tau}^{\rm raw}/\alpha=A_N+B_N$, where $A_N$
contains free evolution on $R$, all outside on-site terms and every outside
interaction disjoint from $R$. Put all seven incident star groups into $B_N$.
Each such star has norm at most $7|\tau|/8$ in the physical/alpha convention, so
\begin{equation}
 \norm{B_N}\le49|\tau|/8,\qquad
 \norm{\alpha_t^\tau(W)-\alpha_t^0(W)}
 \le k|t|,\quad k:=49|\tau|/4.
 \label{eq:duhamel}
\end{equation}
To justify the estimate, differentiate the relative unitary
$e^{it(A_N+B_N)}e^{-itA_N}$ on the common operator domain. The derivative
contains only the bounded perturbation $B_N$; strong integration extends the
identity to every vector and bounds the relative unitary's distance from $I$
by $|t|\norm{B_N}$. Insert bounded $W$ only afterwards to obtain the factor two
in the conjugation bound. This requires no norm derivative of arbitrary local
bounded operators. AQ's already matched unbounded-on-site dynamics theorem
\cite{ns,aq} passes the uniform finite-volume bound to the actual limiting
dynamics. Replacing seven stars by the extensive magnetic norm would lose this
volume uniformity.

The free evolution preserves $R$ and
$c_0(t)=\tfrac14e^{3it}$. Add and subtract the free evolution in the actual
state, then its Haar expectation, keeping the actual mean subtraction:
\begin{align}
 c_\tau(t)&=\omega_\tau(W\alpha_t^\tau(W))-m_\tau^2,\\
 |c_\tau(t)-c_0(t)|&\le k|t|+D+D^2,
 \label{eq:rawcompare}\\
 |c_\tau(t)-c_0(t)|&\le v_\tau+\tfrac14\le B:=\tfrac12+D/2.
 \label{eq:tailbound}
\end{align}
The first bound charges dynamics, local state and centering separately; the
second provides a bounded whole-tail allowance.

\subsection{Poisson transfer and evaluated certificate}
The established Fourier identity \cite{poisson}
\begin{equation}
 \int_{\mathbb R}\frac{s\,e^{itx}}{\pi(s^2+t^2)}\,dt=e^{-s|x|}
 \quad(s>0)
 \label{eq:poisson}
\end{equation}
becomes a heat identity only because the actual ground-centered GNS spectrum
is nonnegative. The positive kernel has mass one and the spectral measure has
finite mass, so spectral Fubini is legitimate. Split the integral at $|t|=L$.
Use \eqref{eq:rawcompare} inside and \eqref{eq:tailbound} on the entire
complement. The two elementary integrals are
\begin{align}
 \int_{-L}^L\frac{s|t|}{\pi(s^2+t^2)}\,dt
 &=\frac{s}{\pi}\log(1+L^2/s^2),\\
 \int_{|t|>L}\frac{s}{\pi(s^2+t^2)}\,dt
 &=\frac2\pi\arctan(s/L)\le\frac{2s}{\pi L}.
\end{align}
Consequently the actual AQ correlation obeys
\begin{equation}
 \boxed{\abs{C_\tau(s)-\tfrac14e^{-3s}}\le
 \mathcal E_\tau(s,L):=
 D+D^2+\frac{ks}{\pi}\log(1+L^2/s^2)
 +B\frac{2s}{\pi L}.}
 \label{eq:certificate}
\end{equation}
Increasing $L$ reduces the tail but increases the logarithmic interior term;
neither may be silently discarded. The formula is a model-specific sufficient
error bound, not a new generic Poisson or Duhamel theorem.

At the frozen AT4 parameters the outward-certified forward error allowance is
approximately $0.000841518704386267$, with the components in
Table~\ref{tab:capcost}. The actual interval is approximately
$[0.0116052483875797,\,0.0132882857963523]$. The exact rational endpoints and
arithmetic proof are in the
\repo{research/round31/forward/at4/output/results.json}{frozen AT4 output}.
Its midpoint radius exceeds $10^{-6}$, so it is not a node satisfying AT3's
input contract. This is failure of this upper certificate; it is not a lower
bound on the true interacting-minus-reference error.

\begin{table}[htbp]
\centering\small
\caption{AT4 forward budget; decimals are readable summaries of outward rational bounds.}
\label{tab:capcost}
\begin{tabular}{lr}
\toprule
Cost & Certified allowance, approximately\\
\midrule
Local state $D$ & $8.08290376865476\times10^{-4}$\\
Mean-square centering $D^2$ & $6.53333333333333\times10^{-7}$\\
Interior dynamics & $7.18276887291769\times10^{-7}$\\
Whole Poisson tail & $3.18567173001654\times10^{-5}$\\
\midrule
Total & $8.41518704386267\times10^{-4}$\\
\bottomrule
\end{tabular}
\end{table}

\subsection{Information and damaging controls}
The inherited abstract measures
$\eta_A=(\delta_2+\delta_4)/8$ and
$\eta_B=\delta_1/32+3\delta_3/16+\delta_5/32$
have common mass, first moment and second moment $(1/4,3/4,5/2)$.
Writing $r=e^{-1}$, their nonzero-time difference is exactly
\begin{equation}
 C_B(1)-C_A(1)=\frac{r(1-r)^4}{32}>0.
\end{equation}
Every deterministic estimator using only those common moments returns the same
number on both. The triangle inequality therefore gives a worst-case error of
at least half their separation, strictly greater than $9/10000$ by directed
exponential arithmetic. This information limit concerns an abstract admissible
moment class, not multiple AQ realizations or nonuniqueness of the actual state.

Centering controls distinguish two operations. If $\widehat m=m+d$, vector
centering adds $d^2$ to the heat correlation, whereas subtracting
$\widehat m^2$ from its raw scalar value adds $-2md-d^2$. For
$m=1/4,d=1/100$, these offsets are $1/10000$ and $-51/10000$, respectively;
no centering leaves the constant vacuum residue $1/16$. Positive energy moments
cannot detect that zero-energy atom. Other exact controls retain all seven
stars, reject a changed physical clock, distinguish nonzero selected vacua from
Haar, retain positive Poisson tail mass, and show that a finite common prefix
does not prove a thermodynamic convergence rate. These fixtures expose invalid
inferences; they are not measured AQ data.

Full derivations:
\repo{research/round31/forward/at4/report.md}{AT4 forward},
\repo{research/round31/reverse/at4/report.md}{AT4 reverse}, and
\repo{research/round31/skeptic/at4.md}{independent model-agent review}.
